\begin{abstract}
As language models increasingly train on web text containing their own outputs, a recursive model$\to$data$\to$model loop arises, yet no method estimates how many such iterations a text has undergone.
We formalize \emph{recursive generation depth estimation} as ordinal classification and construct a benchmark of 180K cell-instances (140K unique passages) spanning three model families, three decoding strategies, and four depths under controlled full-replacement chains at the 1--1.5B parameter scale.
A lightweight pipeline (19 statistical features classified by Random Forest) achieves d1--d3 balanced accuracy of 59--82\% across nine model--strategy cells, with $\mathrm{EDD}_{0.70} = 3$ in 8/9 cells.
Three findings emerge: (1)~depth signal concentrates in perplexity distribution tails, corroborated by classifier convergence and independent cross-perplexity features; (2)~decoding strategy affects cross-condition transfer more than model architecture, suggesting structurally distinct degradation trajectories; (3)~signal attenuates with scale (14B cross-reference accuracy ${\sim}44$\%, $\mathrm{EDD}_{0.70}{=}2$ for nucleus and temperature, 0 for top-$k$), with reference-model choice shifting accuracy by ${\sim}7$\,pp and altering EDD, indicating it is a major confound in cross-scale comparisons.
Our depth estimation claims hold most strongly within-strategy or when the decoding configuration is known; a two-stage routing pipeline recovers $+$11.5\,pp under cross-strategy transfer.
Perturbation-based and membership-inference baselines lack ordinal depth resolution, confirming that depth estimation requires distributional statistics beyond binary detection scores.
\end{abstract}
